\documentclass[conference]{IEEEtran}
\usepackage[utf8]{inputenc}
\usepackage[T1]{fontenc}
\usepackage{cite}
\usepackage{amsmath,amssymb}
\usepackage{graphicx}
\usepackage{booktabs}
\usepackage{url}

\title{Evaluating a Deterministic Validator Plus Single‑Iteration Repair Pipeline for LLM‑Generated Network Configurations}

\author{
\IEEEauthorblockN{Autonomous AI Researcher}
\IEEEauthorblockA{CNSM 2026 Agentic AI Researcher Track}
}

\begin{document}

\maketitle

\begin{abstract}
We preregistered and executed a paired confirmatory experiment (N=300 paired intent\ensuremath{\rightarrow}configuration tasks; preregistration id cand-2026-001-A-prereg-v1, preregistration SHA-256 archived) to evaluate whether a minimal guarded pipeline---single‑shot LLM generation (declared model: gpt-5-mini) followed by a frozen deterministic configuration validator and at most one automated repair iteration---reduces first‑pass misconfiguration relative to plain single‑shot generation. Execution completed 600 episodes and used 301 model calls. Observed outcomes: baseline = 299/300 valid (99.667\%); guarded = 300/300 valid (100.0\%); paired counts n11=299, n10=1, n01=0, n00=0. Exact McNemar two‑sided p = 1.00; absolute paired difference = 0.0033333; paired bootstrap 95\% CI = [0.00, 0.01] (10,000 resamples, seed 1729). The single permitted automated repair corrected the sole invalid baseline output (repair success 1/1; Clopper--Pearson 95\% CI = [0.025, 1.0]). The preregistered planning assumption used for power calculations (planned baseline pass rate = 0.40) was contradicted by the observed near‑ceiling baseline, removing statistical headroom for the preregistered >=50\% relative reduction target. We report preregistered design choices, execution accounting, confirmatory statistics, a post‑hoc inferential sensitivity bound tied to the observed contingency counts, and artifact‑grounded recommendations for future NetOps confirmatory evaluations. All run artifacts and analysis outputs are archived in the execution manifest referenced in the Disclosure Statement.
\end{abstract}

\section{Introduction}
Translating operator intent to device configuration is a core NetOps task. Large language models (LLMs) have been proposed as intent\ensuremath{\rightarrow}configuration translators that can accelerate operations, but probabilistic generation risks misordering, omission, and unsafe commands. Recent literature advocates hybrid patterns that pair probabilistic generators with deterministic validators/simulators and bounded repair loops as operational floor‑safety nets. However, rigorous preregistered evaluations that quantify how tightly bounded guarded pipelines change first‑pass misconfiguration rates under controlled sampling semantics are limited.

We implement and preregister a narrowly scoped paired experiment to evaluate the operational effect of a minimal guarded pipeline composed of: (1) a single shared LLM candidate per task (single‑shot generation), (2) a frozen deterministic validator performing canonicalization and rule/dependency checks, and (3) at most one automated repair attempt when the validator reports violations. The core research question is: Can this minimal guarded pipeline reduce first‑pass misconfiguration rate relative to plain single‑shot generation under a controlled paired design? Our contributions are (i) a preregistered paired experiment with deterministic seed generation and archived per‑task artifacts, (ii) confirmatory statistics derived from those artifacts, (iii) a compact post‑hoc sensitivity quantification tied to the observed contingency counts, and (iv) artifact‑grounded recommendations for future NetOps confirmatory evaluations.

\section{Related Work}
Recent work on LLMs for NetOps and intent‑driven configuration has produced numerous system demonstrations and surveys documenting both promise and concrete failure modes. Empirical demonstrations of intent\ensuremath{\rightarrow}configuration translation report that LLMs can synthesize plausible device configuration templates in lab and testbed settings, yet they frequently surface operationally relevant errors such as command misordering, omission of required safety steps, and overconfident assertions that lack mechanical verification (example system evaluation and discussion: Nguyen et al. 2024 [10.1002/nem.2313]). Complementary survey and review articles consolidate these observations and emphasize the need for grounded verification and governance when applying LLMs to critical network tasks [NIST guidance and surveys, 10.6028/nist.ai.100-2e2023].

Parallel technical threads motivate guarded architectures. Work in adjacent domains (clinical QA, code repair, and general RAG/verification research) shows measurable reductions in factual errors when probabilistic generators are paired with retrieval, deterministic checks, or repair loops; these findings have been ported by several NetOps proposals arguing for generate\ensuremath{\rightarrow}verify\ensuremath{\rightarrow}repair pipelines rather than blind single‑shot deployment [e.g., comparative/procedural studies and framework discussions in the literature pool, 10.36227/techrxiv.176784505.56675782/v1; 10.2139/ssrn.6608518]. Architectural proposals for intent‑aware automated remediation and multi‑agent NetOps orchestration similarly argue that deterministic validators and bounded repair/backout policies are central to safe agentic NetOps workflows (see recent agentic NetOps proposals and intent‑translation demonstrations, e.g., 10.36227/techrxiv.177004277.71795055/v1).

Deterministic network verification provides concrete algorithmic foundations for validators used in guarded pipelines. Prior work on specification‑based and header‑space analyses demonstrates tractable static checks for reachability, access control, and update‑safety properties without requiring live device execution; such algorithms can be integrated to canonicalize generated artifacts into normalized ASTs and to flag violations that trigger repair logic [Beckett et al., 2017 foundation for incremental/specification checks, 10.1145/3098822.3098834]. These deterministic checks differ fundamentally from empirical or emulation tests: they provide provable, rule‑based pass/fail decisions for a defined specification, but they do not by themselves measure runtime control‑plane or packet‑level effects unless combined with higher‑fidelity simulation or emulation.

Limitations in the prior literature motivate the present work. Many demonstrations and proposals are system‑level or exploratory rather than preregistered confirmatory experiments; reported improvements from RAG, verifier‑assisted generation, or agentic architectures often lack paired within‑task statistical controls that isolate the incremental effect of a validator+repair floor on the same generated candidate. Moreover, published comparisons rarely quantify how sampling semantics (shared candidate versus independent draws) and repair budgets affect the number of informative discordant pairs available to paired binary tests. Surveys and architectural reviews therefore call for reproducible evaluation harnesses that combine deterministic validators with controlled sampling semantics and explicit inferential plans [10.6028/nist.ai.100-2e2023].

Our study addresses a tightly scoped portion of this gap by executing a preregistered paired experiment that (a) fixes shared\_initial\_candidate semantics to measure remediation capacity on a single generated plan, (b) uses a frozen deterministic validator to define a binary success criterion tied to intent constraints, and (c) limits repair budget to a single automated iteration. This design directly tests the operational claim that a minimal validator+repair floor reduces first‑pass misconfiguration, and it surfaces the inferential consequences of high baseline pass rates and shared‑candidate pairing for confirmatory McNemar‑style testing.

\section{Methodology}
Preregistration and hypotheses: The experiment was preregistered (cand-2026-001-A-prereg-v1) prior to observing outcomes. Confirmatory hypotheses were: H1 --- the guarded pipeline reduces first‑pass misconfiguration rate by >=50\% (relative) versus single‑shot generation; H2 --- one or fewer automated repair iterations recover >=75\% of initially invalid configurations. The preregistration explicitly used a planning assumption of planned baseline pass rate = 0.40 for power calculations; this numeric assumption is central to the interpretation of the confirmatory outcome and is documented in the preregistration record (see Disclosure).

Design and sampling semantics: We adopted a paired within‑task design with 300 deterministically generated seeds from a frozen task generator. For each seed the same sampled LLM candidate was reused in both arms (shared\_initial\_candidate semantics): baseline (single‑shot generation) and guarded (identical candidate validated and---if invalid---subjected to up to one automated repair attempt). Planned episodes = 600 (two episodes per seed), initial\_generation\_calls\_per\_task = 1, maximum\_repair\_calls\_per\_task = 1. The shared‑candidate pairing was preregistered to isolate the incremental effect of validation and repair on a fixed generated plan; this choice increases within‑pair correlation and reduces the number of informative discordant pairs available to McNemar inference compared with an independent‑draws design.

Model, validator, and scoring: The declared model in the execution contract was gpt-5-mini. The frozen deterministic validator/repair adapter (validator id deterministic\_netops\_validator\_v5) implements rule and dependency checks, canonicalizes configurations to normalized ASTs, and returns binary pass/fail scoring plus violation codes used to trigger and steer the preregistered automated repair adapter. The scoring rule is full intent‑constraint validation: outputs are scored as success only when required command sequences and workflow constraints are satisfied by the validator's canonicalization and rule checks. The preregistered missingness plan allowed one retry for provider failures; primary analysis used the 'complete\_pair\_primary' treatment which drops pairs only if both calls fail. All validator decisions, normalized configurations, and repair traces were archived in the execution manifest and analysis bundle.

Power, estimand, and analysis plan: The primary estimand was the paired\_success\_rate\_difference\_guarded\_minus\_baseline. The preregistered primary test was the exact McNemar two‑sided test (\ensuremath{\alpha}=0.05) on paired pass/fail outcomes; we also preregistered paired bootstrap percentile 95\% CIs (10,000 resamples) for the absolute paired difference. The preregistration included a planning assumption baseline pass rate = 0.40 which informed the sample size and the >=50\% relative reduction target; this assumption is reported here to make the planning calculus transparent. Secondary analyses included exact binomial intervals for repair success, stratified checks, and sensitivity analyses treating failed model calls as failures; all analysis code and per‑task artifacts are archived.

\section{Results}
Execution accounting (artifact‑derived): Planned episodes = 600; completed episodes = 600; complete pairs analyzed = 300. Model calls used = 301 (300 shared initial generations + 1 repair invocation). No provider or infrastructure failures occurred under the preregistered retry policy; no deterministic exclusions for contamination were required under the preregistered checks. Per‑task normalized responses, validator decisions, and repair traces are archived and are the direct source of the numerical counts below.

Primary outcomes (archived counts): Baseline\_success\_count = 299/300 (99.667\%); Guarded\_success\_count = 300/300 (100.0\%). Paired contingency (guarded vs baseline) derived from archived analysis artifacts: n11 = 299 (both success), n10 = 1 (guarded success, baseline fail), n01 = 0 (guarded fail, baseline success), n00 = 0 (both fail). These counts are reproduced in the archived paired contingency artifact used for confirmatory inference.

Inferential results (preregistered): Exact McNemar two‑sided test on the observed contingency gives p = 1.00. The observed absolute paired difference (guarded - baseline success rate) = 1/300 \ensuremath{\approx} 0.0033333. The paired bootstrap percentile 95\% CI for the absolute paired difference = [0.00, 0.01] (10,000 resamples, bootstrap seed = 1729). Repair outcomes (H2): the single allowed automated repair corrected the sole invalid baseline output (repair success 1/1); Clopper--Pearson 95\% CI for 1/1 is [0.025, 1.0], indicating wide uncertainty given the single observed repair.

Post‑hoc sensitivity bound (artifact‑grounded): The total number of discordant pairs observed is n10 + n01 = 1. Under the paired sampling semantics used, this discordant total imposes a hard limit on the resolvable absolute paired difference equal to 1/300 \ensuremath{\approx} 0.00333. Therefore (a) the maximal absolute improvement consistent with the observed discordant count is 1/300; and (b) the preregistered target corresponding to a >=50\% relative reduction from a planned baseline of 0.40 (\ensuremath{\approx} absolute +0.30) would have required on the order of 90 net discordant pairs favoring guarded outputs. This artifact‑grounded bound explains why the preregistered large relative effect was not detectable: the empirical baseline prevalence (299/300) left essentially no statistical headroom for the planned effect size. The archived contingency artifact and bootstrap results support these numerical statements.

Representative diagnostic artifacts: For each episode the archived validator trace contains normalized\_configuration, parsed\_command\_count, required\_command\_count, violation\_codes, and repair\_applied flags. The single baseline failure was traceable in those artifacts and was corrected by invoking the preregistered repair adapter on the shared candidate; the post‑repair canonicalized configuration satisfied the full intent constraints according to the deterministic validator. These representative artifacts are available in the archived execution bundle described in the Disclosure Statement.

\section{Discussion}
Interpretation and operational implications: The guarded pipeline corrected the single observed invalid baseline output, demonstrating that a frozen deterministic validator plus a bounded automated repair iteration can remediate validator‑detected violations in at least some instances. However, because the empirical baseline pass rate was far higher than the preregistered planning assumption, the run does not provide confirmatory evidence for the preregistered >=50\% relative reduction. Practically, this implies that when baseline error rates are uncertain, confirmatory designs must focus on producing informative discordant pairs rather than simply increasing N under optimistic failure assumptions.

Sampling semantics and pairing tradeoffs: The preregistered shared\_initial\_candidate design intentionally reduced cross‑arm variability to isolate validator/repair effects on a single generated plan. That choice is operationally motivated (real workflows often assess and remediate a single proposed plan) but reduces the number of statistically informative discordant pairs. Independent draws per arm would increase discordance and therefore statistical information for McNemar tests but at the cost of within‑pair heterogeneity. Future confirmatory studies should align pairing semantics with the targeted operational claim and incorporate their inferential consequences into power calculations.

Repair budget and conditional estimands: We treated repair effectiveness as H2 with a single allowed automated repair iteration. The observed repair success (1/1) is artifact‑grounded but statistically imprecise. When repair effectiveness is a principal operational claim, we recommend treating repair success as a conditional estimand (conditional on initial failure) and powering the study to observe a sufficient number of initial failures to estimate that conditional probability with narrow confidence intervals.

Threats to validity and externality: Internal validity is affected by the preregistered shared\_candidate semantics and the deterministic replacement rules in the preregistration. Construct validity is limited because the deterministic validator performs static canonicalization and rule checks rather than executing configurations in live devices; dynamic control‑plane or packet‑level consequences were not measured. External validity is limited by the synthetic task bank, constrained vendor template space, and a single model family (gpt-5-mini) and validator implementation. Conclusion validity is constrained by the paucity of informative discordant pairs (one), which leaves large preregistered effects unverifiable under the observed data. We document these threats and their artifact bases in the archived manifest.

Design recommendations for future confirmatory NetOps evaluations (artifact‑grounded): (1) Calibrate task banks to target a modest initial‑failure prevalence or a target discordant‑pair count aligned with the planned effect size; concrete options include increasing task difficulty, adversarial augmentation, or stratified sampling by difficulty. (2) Explicitly preregister pairing semantics (shared candidate versus independent draws) and repair budgets and include their expected effect on discordant counts in power computations. (3) When operational claims require runtime semantic assurances, incorporate network emulation or simulation that exercises packet‑level and control‑plane consequences instead of relying solely on static validators. (4) Archive per‑task validator traces and repair prompts systematically to enable post‑hoc diagnostic analyses and reproducibility.

\section{Limitations}
\begin{itemize}
\item Synthetic task bank and constrained vendor‑template scope limit external validity to broad production environments.
\item Single LLM (gpt-5-mini) and a single validator implementation restrict generality across model families and validator designs.
\item The preregistered power plan assumed a baseline pass rate = 0.40; the observed baseline (299/300) contradicts that assumption and removed headroom to detect the preregistered >=50\% relative reduction.
\item Sampling semantics (shared\_initial\_candidate) reused the same initial candidate across paired arms, increasing within‑pair correlation and reducing discordant pairs available to McNemar inference.
\item Validation used deterministic configuration/constraint checks rather than full dynamic emulation of runtime semantic effects; actual runtime consequences of configurations were not executed and remain unmeasured.
\item H2 repair‑success evidence is based on a single initially invalid case (1/1 \ensuremath{\rightarrow} 100\%), yielding a wide Clopper--Pearson 95\% CI = [0.025, 1.0]; larger counts of initial failures are required to estimate repair effectiveness precisely.
\end{itemize}

\section{Conclusion}
We executed a preregistered paired experiment (N=300 pairs) comparing single‑shot LLM generation (gpt-5-mini) to a guarded pipeline consisting of a frozen deterministic validator and a single automated repair iteration. Baseline single‑shot generation yielded 299/300 valid outputs and the guarded pipeline yielded 300/300 valid outputs. The single automated repair corrected the lone invalid baseline output, but the observed near‑ceiling baseline contradicted the preregistered planning assumption (baseline pass rate = 0.40) and removed statistical headroom to detect the preregistered >=50\% relative reduction. Future confirmatory NetOps evaluations should (a) calibrate task difficulty to produce sufficient initial failures or target discordant counts, (b) preregister pairing semantics and repair budgets explicitly, and (c) include emulation to measure runtime semantic effects when operational deployment claims are intended. All raw outputs, validator traces, repair traces, the execution manifest, and analysis artifacts are archived and enumerated in the Disclosure Statement to permit reproduction of the reported counts and intervals.

\section*{Disclosure Statement}
This study was executed autonomously after the autonomy lock under the archived master prompt (provenance/master\_prompt.txt; SHA-256: 1872df1e\allowbreak{}1805d2d9\allowbreak{}6940456c\allowbreak{}a016bd66\allowbreak{}5d1d5196\allowbreak{}add77f5a\allowbreak{}cdf1582b\allowbreak{}b39b15ba). Human role: the Research Director selected the topic family (Generative AI and LLMs for NetOps) and approved the master prompt prior to the autonomy lock; no manual human editing of manuscript technical content, figures, tables, or references occurred after the lock. Preregistration: cand-2026-001-A-prereg-v1 (the preregistration record was created before observing final outcomes; preregistration SHA-256 is recorded in the execution manifest). Declared model and tooling: gpt-5-mini (declared\_model\_version), adapter family hosted\_netops\_gvr\_v1, frozen deterministic validator/repair adapter deterministic\_netops\_validator\_v5. Execution accounting (archived): planned episodes = 600, completed episodes = 600, complete pairs = 300, model\_calls\_used = 301 (300 shared initial generation calls + 1 repair invocation), repair-stage invocations = 1. The execution manifest and analysis artifact bundle enumerate all raw model responses, per‑task validator traces (normalized configurations, parsed\_command\_count, required\_command\_count, violation\_codes, repair\_applied flags), repair prompts/responses, execution logs, and analysis artifacts. The archived execution manifest and analysis bundle are the authoritative source for the numeric counts reported in this manuscript. The autonomous pipeline performed literature search, preregistration, execution, analysis, AI‑peer‑review style checks, and manuscript drafting autonomously after the autonomy lock in accordance with the CNSM Agentic AI Researcher track requirements. The preregistered planning assumption used in power calculations (planned baseline pass rate = 0.40) is explicitly reported here because it materially affects interpretation of the confirmatory result.

\begin{thebibliography}{99}
\bibitem{ref1} Nguyen Tu, Sukhyun Nam, James Won‐Ki Hong, "Intent‐Based Network Configuration Using Large Language Models", 2024, 10.1002/nem.2313
\bibitem{ref2} Oghenekeno Hilkiah Okula, ""The Era of AI Agents for Network Engineering": Intent-Aware Auto Remediation for Network Disruption or Failures Using AI Agents, Large Language Models (LLM) and Model Context Protocol (MCP) Mechanisms", 2026, 10.36227/techrxiv.177004277.71795055/v1
\bibitem{ref3} Kunal Dhanda, "MedRAG: Retrieval-Augmented Generation for Medical QA-Comparing Base and RAG-Augmented LLM Performance on Evidence from Peer-Reviewed Clinical Research", 2026, 10.2139/ssrn.6608518
\bibitem{ref4} Goutam Adwant, Manjari Srivastav, "Evidence Coverage Evaluator: A Comprehensive Framework for Assessing Grounding Quality in Retrieval-Augmented Generation Systems", 2026, 10.36227/techrxiv.176784505.56675782/v1
\bibitem{ref5} Ryan Beckett, Aarti Gupta, Ratul Mahajan, David Walker, "A General Approach to Network Configuration Verification", 2017, 10.1145/3098822.3098834
\bibitem{ref6} Apostol Vassilev, Alina Oprea, Alie Jean Fordyce, Hyrum S. Anderson, "Adversarial machine learning :", 2024, 10.6028/nist.ai.100-2e2023
\end{thebibliography}

\end{document}
